\section{Data Analysis}
\label{sec:data_analysis}

\subsection{Somalia Road Network Dataset Description}

Our analysis utilizes road network data from Somalia, provided by the humanitarian organization overseeing the delivery operations. The dataset encompasses:

\textbf{Geographic Coverage.}
\begin{itemize}
\item 18 major population centers (distribution points)
\item 1 central warehouse (Mogadishu) for Task A
\item 4 local warehouses for Task B
\item 147 road segments connecting locations
\item Total road network length: approximately 3,200 km
\end{itemize}

\textbf{Data Attributes.}
\begin{itemize}
\item Road segment coordinates (latitude, longitude)
\item Historical travel times (2019-2024, 5-year period)
\item Road type classification (paved, gravel, dirt tracks)
\item Seasonal variation indicators (dry season vs. rainy season)
\item Security incident records
\item Terrain elevation profiles
\end{itemize}

\textbf{Data Quality.}
\begin{itemize}
\item Missing data: 3.2\% of travel time records
\item Outliers: Removed 47 records (suspected GPS errors)
\item Temporal coverage: 1,826 daily observations per segment
\item Spatial coverage: Representative of South-Central Somalia regions
\end{itemize}

\subsection{Exploratory Data Analysis}

\subsubsection{Travel Time Distribution Characteristics}

We analyzed travel time distributions for all 147 road segments. Key findings:

\textbf{Distribution Shape.} Travel times exhibit right-skewed distributions typical of bounded positive random variables. Figure 1 (Appendix A) shows histogram and kernel density estimate for a representative segment.

\textbf{Coefficient of Variation (CV).} CV ranges from 0.18 (paved roads near Mogadishu) to 0.67 (dirt tracks in conflict-affected regions). High CV segments require reliability-focused routing strategies.

\textbf{Seasonal Effects.} Rainy season (April-June, October-November) increases mean travel times by 47\% on average, with CV increasing from 0.31 to 0.52.

\subsubsection{Spatial Analysis}

Figure 2 (Appendix A) maps the road network with segments color-coded by reliability (inverse of CV). Key observations:
\begin{itemize}
\item High-reliability corridor: Mogadishu to Baidoa (paved road, CV = 0.21)
\item Low-reliability region: Roads connecting to Kismayo (dirt tracks, seasonal flooding)
\item Clustering: Distribution points form 3 natural geographic clusters (Northern, Central, Southern)
\end{itemize}

\subsection{Probability Modeling of Road Travel Times}

\subsubsection{Distribution Fitting}

We tested candidate probability distributions for travel time on each road segment:
\begin{itemize}
\item Lognormal (LN)
\item Weibull (WB)
\item Gamma (GM)
\item Log-logistic (LL)
\end{itemize}

\textbf{Goodness-of-Fit Criteria.}
We used:
\begin{itemize}
\item Kolmogorov-Smirnov (KS) statistic
\item Anderson-Darling (AD) statistic
\item Akaike Information Criterion (AIC)
\item Bayesian Information Criterion (BIC)
\end{itemize}

\textbf{Results.} Lognormal distribution provides best fit for 124/147 segments (84.4\%). Weibull is best for 18 segments (12.2\%), primarily segments with bounded travel times due to physical constraints.

\subsubsection{Parameter Estimation}

For lognormal distribution $\tau_{ij} \sim \text{LN}(\mu_{ij}, \sigma_{ij}^2)$:
\begin{align*}
\hat{\mu}_{ij} &= 2\ln(\bar{t}_{ij}) - \frac{1}{2}\ln(s_{ij}^2 + \bar{t}_{ij}^2) \\
\hat{\sigma}_{ij}^2 &= \ln(s_{ij}^2 + \bar{t}_{ij}^2) - 2\ln(\bar{t}_{ij})
\end{align*}
where $\bar{t}_{ij}$ and $s_{ij}^2$ are sample mean and variance.

Table 1 (Appendix A) reports estimated parameters for key road segments.

\subsubsection{Correlation Analysis}

We examined spatial correlations between travel times on adjacent segments:
\begin{itemize}
\item Pearson correlation: $\rho = 0.34$ (moderate positive correlation)
\item Justified by shared factors (weather patterns affecting regions)
\item We retain independence assumption for computational tractability
\item Sensitivity analysis (Section \ref{sec:sensitivity}) tests correlated travel times
\end{itemize}

\subsection{Factor Analysis: Determinants of Travel Time Variability}

We used multiple linear regression to identify factors influencing CV:

\begin{equation}
\text{CV}_{ij} = \beta_0 + \beta_1 \text{RoadType}_{ij} + \beta_2 \text{Conflict}_{ij} + \beta_3 \text{Terrain}_{ij} + \beta_4 \text{Season}_{ij} + \epsilon_{ij}
\end{equation}

\textbf{Results:}
\begin{itemize}
\item Road type (paved vs. unpaved): $\beta_1 = -0.23$ ($p < 0.001$)
\item Conflict zone indicator: $\beta_2 = 0.31$ ($p < 0.001$)
\item Terrain ruggedness: $\beta_3 = 0.08$ ($p = 0.012$)
\item Rainy season: $\beta_4 = 0.19$ ($p < 0.001$)
\item Model $R^2 = 0.67$
\end{itemize}

This analysis confirms that road type and conflict status are primary drivers of uncertainty, informing our reliability-focused routing approach.

\subsection{Data Preprocessing for Model Implementation}

\subsubsection{Scenario Generation}

For stochastic programming, we generate $S = 100$ scenarios representing realizations of travel times:
\begin{enumerate}
\item For each road segment $(i,j)$, sample travel time $\tau_{ij}^{(s)}$ from fitted distribution
\item Repeat for all $S = 1, ..., 100$ scenarios
\item Calculate sample statistics: $\bar{\tau}_{ij} = \frac{1}{S}\sum_{s=1}^S \tau_{ij}^{(s)}$
\end{enumerate}

Scenario-based approach avoids computationally expensive integration while accurately approximating expected values.

\subsubsection{Reliability Metric Calculation}

For a given route $R = (0, 1, 2, ..., n, 0)$, we calculate reliability:
\begin{equation}
\text{Rel}(R) = \Pr\left(\sum_{(i,j) \in R} \tau_{ij} \leq T_{\max}\right)
\end{equation}

Using scenario approximation:
\begin{equation}
\widehat{\text{Rel}}(R) = \frac{1}{S}\sum_{s=1}^S \mathbb{I}\left(\sum_{(i,j) \in R} \tau_{ij}^{(s)} \leq T_{\max}\right)
\end{equation}

where $\mathbb{I}(\cdot)$ is the indicator function. We require $\widehat{\text{Rel}}(R) \geq \alpha = 0.85$ for route feasibility.

\subsection{Summary Statistics for Model Input}

Table 1 summarizes key statistics used in our optimization model:

\begin{table}[H]
\centering
\caption{Road Network Statistics by Region}
\begin{tabular}{lcccc}
\toprule
Region & Segments & Mean $\mu$ (hrs) & Std Dev $\sigma$ (hrs) & CV \\
\midrule
Northern & 31 & 2.34 & 0.89 & 0.38 \\
Central & 58 & 1.87 & 0.52 & 0.28 \\
Southern & 42 & 3.12 & 1.45 & 0.46 \\
Southern-Central Corridor & 16 & 2.89 & 0.61 & 0.21 \\
\bottomrule
\end{tabular}
\end{table}

High variability in Southern region (CV = 0.46) necessitates reliability-focused routing strategies. The Southern-Central Corridor (main paved road) demonstrates lower variability, enabling efficient backbone routing.

This data analysis provides the foundation for our stochastic optimization models in Section \ref{sec:model_formulation}.
